\begin{ouverture}[calcul]{4}{Carrés et racines carrées}{Partie 1 · Nombres et calculs}
{Carrés parfaits · Valeur exacte · Encadrement}
\objectif{Connaître les carrés des entiers de 1 à 12.}
\objectif{Définir et encadrer une racine carrée.}
\objectif{Passer d’une aire à une longueur sans arrondir trop tôt.}
\prerequis{Multiplication et aire d’un carré.}
\end{ouverture}
\pagestyle{otMain}\essentiel\fontsize{11}{15}\selectfont
\phantomsection\label{cours-4-1}
\essSection{1}{Du carré à sa racine}
\begin{definition}Le carré d’un nombre $a$ est $a^2=a\times a$. Il est toujours positif ou nul. Pour un nombre $b\ge0$, $\sqrt b$ désigne l’unique nombre positif ou nul dont le carré vaut $b$. Ainsi $\sqrt{49}=7$ et $\sqrt0=0$.\end{definition}
\begin{center}\begin{tabular}{c|rrrrrr}\toprule
$n$ & 1 & 2 & 3 & 4 & 5 & 6\\\midrule
$n^2$ & 1 & 4 & 9 & 16 & 25 & 36\\\midrule
$n$ & 7 & 8 & 9 & 10 & 11 & 12\\\midrule
$n^2$ & 49 & 64 & 81 & 100 & 121 & 144\\\bottomrule
\end{tabular}\end{center}
\begin{exemple}Un carré d’aire $30\ \mathrm{cm}^2$ a un côté de $\sqrt{30}$ cm. Ce nombre est positif ; on ne choisit pas une longueur négative.\end{exemple}
\begin{attention}$\sqrt{25}=5$ seulement. Les deux nombres 5 et $-5$ ont pour carré 25, mais le symbole racine carrée désigne le nombre non négatif. Aucun nombre réel n’a pour carré $-9$.\end{attention}
\clearpage
\phantomsection\label{cours-4-2}
\essSection{2}{Encadrer et approcher}
\begin{savoirfaire}[Utiliser les carrés connus]Comme $25<30<36$, on sait que $5<\sqrt{30}<6$. Pour affiner, comparer des carrés décimaux : $5{,}4^2=29{,}16$ et $5{,}5^2=30{,}25$, donc $5{,}4<\sqrt{30}<5{,}5$.\end{savoirfaire}
\begin{exemple}La calculatrice donne $\sqrt{30}\approx5{,}477225\ldots$. Au centième, on écrit $\sqrt{30}\approx5{,}48$. Le signe « $=$ » reste réservé à l’égalité exacte. Le périmètre exact du carré est $4\sqrt{30}$ cm, environ 21,91 cm.\end{exemple}
\begin{savoirfaire}[Relier aire et longueur]Calculer d’abord l’aire cherchée si nécessaire, puis prendre sa racine carrée. Dans une formule de longueur, vérifier que l’expression sous la racine représente bien une aire positive ou nulle. Garder la valeur exacte dans les calculs intermédiaires.\end{savoirfaire}
\begin{attention}Une racine ne se répartit pas sur une somme : $\sqrt{9+16}=5$, alors que $\sqrt9+\sqrt{16}=7$. Les transformations algébriques générales des racines ne sont pas un objectif de cette édition.\end{attention}

\clearpage\rubrique{Découverte}
\begin{blocLivre}{ex-4-1}\begin{exercice}[1]{Les carrés connus}
Calculer $7^2$, $11^2$ et $12^2$.
\end{exercice}\end{blocLivre}
\begin{blocLivre}{ex-4-2}\begin{exercice}[1]{Lire trois racines}
Donner $\sqrt{64}$, $\sqrt{100}$ et $\sqrt{0}$.
\end{exercice}\end{blocLivre}
\begin{blocLivre}{ex-4-3}\begin{exercice}[1]{Un signe à expliquer}
Pourquoi $\sqrt{81}$ ne vaut-il pas $-9$ ?
\end{exercice}\end{blocLivre}
\begin{blocLivre}{ex-4-4}\begin{exercice}[1]{Encadrer entre deux entiers}
Encadrer $\sqrt{53}$ entre deux entiers consécutifs.
\end{exercice}\end{blocLivre}
\begin{blocLivre}{ex-4-5}\begin{exercice}[1]{De l’aire au côté}
Un carré a une aire de $121\ \mathrm{cm}^2$. Quelle est la longueur de son côté ?
\label{fig-ex-4-5}\figureGeo{carre-121}
\end{exercice}\end{blocLivre}

\rubrique{Entraînement}
\begin{blocLivre}{ex-4-6}\begin{exercice}[2]{Une valeur approchée}
Donner une valeur approchée de $\sqrt{19}$ au centième et un encadrement par deux dixièmes consécutifs.
\end{exercice}\end{blocLivre}
\begin{blocLivre}{ex-4-7}\begin{exercice}[2]{Périmètre exact}
Un carré a une aire de $18\ \mathrm{m}^2$. Donner son périmètre exact, puis arrondi au centième de mètre.
\label{fig-ex-4-7}\figureGeo{carre-18}
\end{exercice}\end{blocLivre}
\begin{blocLivre}{ex-4-8}\begin{exercice}[2]{Deux aires}
Comparer les côtés de deux carrés d’aires 42 et 45 cm$^2$, puis encadrer leur différence entre 0 et 1 cm.
\end{exercice}\end{blocLivre}
\begin{blocLivre}{ex-4-9}\begin{exercice}[2]{Détecter une fausse règle}
Camille propose $\sqrt{36+64}=\sqrt{36}+\sqrt{64}$. Calculer les deux membres et conclure.
\end{exercice}\end{blocLivre}
\begin{blocLivre}{ex-4-10}\begin{exercice}[2]{Arrondi ou troncature}
Un élève écrit $\sqrt{50}=7{,}07$. Corriger l’écriture et expliquer si 7,07 est un arrondi au centième.
\end{exercice}\end{blocLivre}

\rubrique{Problèmes}
\begin{blocLivre}{ex-4-11}\begin{exercice}[3]{Un jardin carré}
Un jardin carré de 72 m$^2$ est entouré d’une clôture sauf sur une ouverture de 2 m. Quelle longueur de clôture prévoir au dixième de mètre ?
\label{fig-ex-4-11}\figureGeo{jardin-72}
\end{exercice}\end{blocLivre}
\begin{blocLivre}{ex-4-12}\begin{exercice}[3]{Passer par une conversion}
Un carré a une aire de $6400\ \mathrm{cm}^2$. Donner son côté en centimètres, puis en mètres.
\end{exercice}\end{blocLivre}
\begin{blocLivre}{ex-4-13}\begin{exercice}[3]{Une aire multipliée}
Un carré de côté 3 cm est remplacé par un carré d’aire cinq fois plus grande. Donner le nouveau côté exact et approché au millimètre.
\label{fig-ex-4-13}\figureGeo{carre-45}
\end{exercice}\end{blocLivre}
\begin{blocLivre}{ex-4-14}\begin{exercice}[3]{Un encadrement précis}
Sans touche racine, encadrer $\sqrt{10}$ entre deux centièmes consécutifs en testant 3,16 et 3,17.
\end{exercice}\end{blocLivre}
\begin{blocLivre}{ex-4-15}\begin{exercice}[3]{Deux projets}
Une plaque carrée de 50 cm$^2$ doit entrer dans une boîte dont le fond intérieur est un carré de côté 7 cm, les côtés restant parallèles. Est-ce possible ? Quelle aire maximale conviendrait ?
\label{fig-ex-4-15}\figureGeo{plaque-50}
\end{exercice}\end{blocLivre}

\clearpage
